\documentclass[a4paper,12pt]{article}

% Кодировка и язык
\usepackage[utf8]{inputenc}
\usepackage[T2A]{fontenc}
\usepackage[russian]{babel}

% Геометрия страницы
\usepackage{geometry}
\geometry{left=2.5cm,right=2cm,top=2cm,bottom=2cm}

% Графика и таблицы
\usepackage{graphicx}
\usepackage{booktabs}
\usepackage{array}

% Листинги кода
\usepackage{listings}
\usepackage{xcolor}

% Гиперссылки
\usepackage{hyperref}
\hypersetup{
    colorlinks=true,
    linkcolor=blue,
    urlcolor=blue,
}

% Настройки листингов
\lstset{
    basicstyle=\ttfamily\small,
    keywordstyle=\color{blue}\bfseries,
    commentstyle=\color{green!60!black}\itshape,
    stringstyle=\color{red},
    numbers=left,
    numberstyle=\tiny\color{gray},
    breaklines=true,
    frame=single,
    language=Python,
    showstringspaces=false,
    tabsize=4,
}

% Заголовок
\title{Лабораторная работа №2\\Документирование кода на Python}
\author{ИСУ: 561636}
\date{\today}

\begin{document}

\maketitle
\thispagestyle{empty}

\newpage
\tableofcontents
\newpage

\section{Цель работы}

Изучить принципы документирования кода на языке Python. Научиться
писать docstring'и к функциям, использовать осмысленные имена
переменных и функций, оформлять проектную документацию в формате
Markdown и верстать отчёты в системе LaTeX.

\section{Ход работы}

\subsection{Анализ исходного файла}

Исходный файл \texttt{mysterious\_file.py} содержал функции с
неинформативными именами (\texttt{h}, \texttt{k}, \texttt{z},
\texttt{f}, \texttt{s}) и без документации. Для каждой функции
было выполнено:

\begin{enumerate}
    \item Прочитан код и определено назначение функции.
    \item Определены типы параметров и возвращаемого значения.
    \item Переименованы функция и переменные в осмысленные имена.
    \item Добавлен docstring с описанием, параметрами, возвращаемым
          значением и примером вызова.
\end{enumerate}

\subsection{Пример документированной функции}

Ниже приведён пример одной из документированных функций.

\begin{lstlisting}[caption={Функция \texttt{is\_even}}]
def is_even(number):
    """
    Checks if a number is even.

    Args:
        number (int): Number to check.

    Returns:
        bool: True if the number is even, otherwise False.

    Example:
        >>> is_even(4)
        True
        >>> is_even(7)
        False
    """
    return number % 2 == 0
\end{lstlisting}

\subsection{Список переименованных функций}

В таблице~\ref{tab:rename} приведено соответствие старых и новых имён.

\begin{table}[h]
\centering
\begin{tabular}{|l|l|l|}
\hline
\textbf{Старое имя} & \textbf{Новое имя} & \textbf{Назначение} \\
\hline
\texttt{buy\_products} & \texttt{buy\_products} & Расчёт сдачи \\
\texttt{check} & \texttt{is\_weekend} & Проверка выходного \\
\texttt{h} & \texttt{is\_even} & Проверка чётности \\
\texttt{z} & \texttt{find\_multiples\_of\_three} & Поиск кратных 3 \\
\texttt{sentiment\_score} & \texttt{calculate\_sentiment\_score} & Анализ тональности \\
\texttt{f} & \texttt{to\_binary\_string} & Перевод в двоичную \\
\hline
\end{tabular}
\caption{Переименование функций}
\label{tab:rename}
\end{table}

\subsection{Создание документации}

В каталоге \texttt{docs} создан файл \texttt{functions.md} в формате
Markdown. Документация включает:

\begin{itemize}
    \item описание проекта;
    \item требования к системе;
    \item структуру проекта;
    \item описание всех функций с примерами использования;
    \item историю изменений с хешами коммитов.
\end{itemize}

\subsection{История коммитов}

В таблице~\ref{tab:commits} приведена история коммитов.

\begin{table}[h]
\centering
\begin{tabular}{|c|l|l|}
\hline
\textbf{№} & \textbf{Хеш коммита} & \textbf{Описание} \\
\hline
1 & \texttt{335f9ce} & All the functions in mysterious\_file.py are documented \\
2 & \texttt{c3e4551} & Documentation for all functions has been added \\
\hline
\end{tabular}
\caption{История коммитов}
\label{tab:commits}
\end{table}

\section{Вывод}

В ходе лабораторной работы были изучены принципы документирования
кода: написание docstring'ов, использование осмысленных имён,
оформление проектной документации в Markdown и верстка отчётов
в LaTeX. Все функции файла \texttt{mysterious\_file.py} были
задокументированы, создана документация в каталоге \texttt{docs},
оформлен отчёт в LaTeX.

\end{document}